\begin{prob}[(3 points)]

    Consider the following code, which takes in a list of integers and moves
    all of the odd numbers to the front and all of the even numbers to the end,
    in-place.

    Remember that the \verb!%! operator computes the remainder of division.
    For example \verb!5 % 2 == 1! and \verb!4 % 2 == 0!. This can be
    used to check if a number is even; if \verb!x % 2 == 0!, then
    \python{x} is even.

    \inputminted[lastline=18]{python}{\thisdir/code.py}

    \begin{subprobset}

        \begin{subprob}(2 points)
            Which of the following are valid loop invariants for the \python{while}-loop? Select all that apply.

            \newcommand{\choicea}{\choice  After the $\alpha$th iteration, the first $\alpha$ elements of \texttt{numbers} are odd.}
            \newcommand{\choiceb}{\choice  After the $\alpha$th iteration, the last $\alpha$ elements of \texttt{numbers} are even.}
            \newcommand{\choicec}{\correctchoice After each iteration, the elements in \python{numbers[:odd_barrier]} are odd.}
            \newcommand{\choiced}{\choice After each iteration, the elements in \python{numbers[odd_barrier+1:]} are even.}
            \newcommand{\choicee}{\choice After each iteration, the elements in \python{numbers[:even_barrier]} are odd.}
            \newcommand{\choicef}{\correctchoice After each iteration, the elements in \python{numbers[even_barrier+1:]} are even.}
            \newcommand{\choiceg}{\choice After the $\alpha$th iteration, \python{even_barrier - odd_barrier = } $\alpha$.}
            \newcommand{\choiceh}{\correctchoice After the $\alpha$th iteration, \python{even_barrier - odd_barrier = n - 1 - } $\alpha$.}

            \onlyversion{A}{
                \begin{choices}[rectangle]
                    \choicea
                    \choiceb
                    \choicec
                    \choiced
                    \choicee
                    \choicef
                    \choiceg
                    \choiceh
                \end{choices}
            }

            \onlyversion{B}{
                \begin{choices}[rectangle]
                    \choiceb
                    \choicea
                    \choiced
                    \choicec
                    \choicef
                    \choicee
                    \choiceh
                    \choiceg
                \end{choices}
            }

        \end{subprob}

        \begin{subprob}
            How many iterations will the \python{while}-loop in
            \python{separate_odd_from_even} make if \python{numbers} has length
            $100$? Your answer should be a number.

            \inlineresponsebox{$99$}
        \end{subprob}

    \end{subprobset}

\end{prob}
